\documentclass[11pt]{article}
\usepackage[margin=25mm]{geometry}
\usepackage{fontspec}
\setmainfont{Latin Modern Roman}
\usepackage{amsmath,amssymb,mathtools}
\usepackage{unicode-math}
\setmathfont{Latin Modern Math}
\usepackage{xurl,hyperref}
\hypersetup{colorlinks=true,urlcolor=blue}
\setcounter{secnumdepth}{0}
\setlength{\parindent}{0pt}
\setlength{\parskip}{5pt}
\setlength{\emergencystretch}{3em}
\newcommand{\sref}[2]{\hyperlink{source:#1:#2}{[S#2]}}
\newcommand{\cataloguescope}[1]{\paragraph{Recorded target.}#1}
\newcommand{\R}{\mathbb R}
\newcommand{\C}{\mathbb C}
\newcommand{\N}{\mathbb N}
\newcommand{\Z}{\mathbb Z}
\newcommand{\Q}{\mathbb Q}
\newcommand{\F}{\mathbb F}
\newcommand{\eps}{\varepsilon}
\begin{document}
% UnsolvedMath ID 30006043; source catalogue code OWR-14298591-007
\setcounter{section}{30006042}
\section{Dynamics Near Generic Spacetime Singularities}\label{30006043}
{\small UnsolvedMath ID 30006043 · Mathematical Physics}
\subsection{Definitions and mathematical statement}

\paragraph{Shared notation.}$\mathbb N=\{1,2,\ldots\}$; the standard number fields have their usual meanings. All additional conventions are stated locally.


The Belinski--Khalatnikov--Lifshitz picture describes local near-singularity evolution as successive approximately homogeneous Kasner epochs with transitions. Generic Einstein solutions may have spatial dependence and need not satisfy a symmetry reduction. A rigorous dynamical description must specify asymptotic variables, convergence and the class of singularities.

\paragraph{Catalogue wording retained for traceability.}
Give a rigorous dynamical description of generic solutions of the Einstein equations near spacetime singularities, including the extent to which Belinski–Khalatnikov–Lifshitz dynamics governs their asymptotic behavior.

\paragraph{Scope and source qualification.}This is a programme covering the extent and limitations of a heuristic, not a claim that every singularity obeys one prescribed BKL sequence.
\subsection{Short English statement}
Determine how well BKL dynamics describes generic Einstein solutions approaching spacetime singularities.
\subsection{Research attempt}
\paragraph{Model, date and investigated scope.}
OpenAI GPT-6 Astra (configured model \texttt{gpt-6-astra}), ultra reasoning; generated 11 October 2026. Derived the vacuum Kasner constraints and checked true versus coordinate singularities in explicit homogeneous metrics.

\paragraph{Sources and literature check.}
Read the full supplied definition, catalogue wording, source qualifications and listed references for 30006043. No fresh status or novelty search was performed. The argument below proves its auxiliary statements directly and does not import an unverified external theorem.

\paragraph{A homogeneous vacuum consistency calculation.} Consider the diagonal metric \(g=-dt^2+\sum_{i=1}^3t^{2p_i}dx_i^2\), \(t>0\), and write \(S=\sum_i p_i\), \(Q=\sum_i p_i^2\). Its nonzero basic connection coefficients are \(\Gamma^t_{ii}=p_it^{2p_i-1}\) and \(\Gamma^i_{ti}=p_i/t\). Substitution into the Ricci formula gives
\[
 R_{tt}=\frac{S-Q}{t^2},\qquad
 R_{ii}=p_i(S-1)t^{2p_i-2},\qquad R_{ij}=0\ (i\ne j).
\]
Thus the Kasner conditions \(S=Q=1\) give a vacuum solution. If at least one \(p_i\ne0\), the same equations force those conditions; the separate all-zero case is static flat spacetime.
\paragraph{Curvature countercheck.} In an orthonormal frame the squared curvature is
\[
 R_{abcd}R^{abcd}=\frac4{t^4}\left(\sum_i p_i^2(p_i-1)^2+\sum_{i<j}p_i^2p_j^2\right).
\]
The bracket vanishes for permutations of \((1,0,0)\), which are flat coordinate cases despite the degenerating metric coefficient. For \((-1/3,2/3,2/3)\) it is positive, proving a curvature singularity at \(t=0\). These examples show why coordinate degeneration and curvature blowup must be distinguished.
\paragraph{BKL gap.} This family has no spatial derivative terms and no transitions between Kasner epochs. It therefore cannot establish generic BKL behavior. The first missing estimate is control of spatially inhomogeneous Einstein terms relative to velocity-dominated terms, including possible transition and exceptional regions. The source asks the extent and limitations of BKL asymptotics for a specified generic class, not merely existence of singular homogeneous examples.

\paragraph{Verification and adversarial audit.}
Checked Ricci coefficients, the all-zero exception, Kasner constraints and flat versus nonflat curvature invariants; no genericity claim. Model self-audit only; no independent human mathematical verification. Applied contribution prompt A before drafting and B afterwards. Checked Ricci coefficients, the all-zero exception, Kasner constraints and flat versus nonflat curvature invariants; no genericity claim. First remaining gap: Control spatial derivative terms and transitions in a precise generic class of inhomogeneous Einstein solutions near singularity. The inherited source status is preserved; original-target completion remains zero.

Coordinator checked the connection, Ricci components and generic versus flat Kasner curvature cases. This additional check was model-assisted, not independent human verification.

\paragraph{Tools, process and contributions.}
Python 3 and local file tools used to read definitions and serialize manuscripts. Calculations are symbolic reasoning; no numerical experiment or theorem prover. Contribution prompt A was followed by prompt B. GPT-6 Astra with ultra reasoning drafted the research and performed its stated self-audit; any additional coordinator contribution is identified in the verification and audit. A GPT-6-based Codex coordinator assembled the manuscript and checked repository integration; its exact variant and reasoning configuration were not exposed. No human mathematical contribution, independent human proof check or formal proof certification is supplied. Complete derivations and the consulted source references are included here; private reasoning transcripts are not published.

\paragraph{Final status and first remaining gap.}
\textbf{UNRESOLVED}. Control spatial derivative terms and transitions in a precise generic class of inhomogeneous Einstein solutions near singularity.

\section{Completion Estimate}
\noindent\textbf{COMPLETION: 0\%}
This is an elementary auxiliary calculation or an explicit test of a restricted strategy. No new portion of the full source target is established.

\subsection{Sources}
\begin{itemize}
\item[\hypertarget{source:30006043:1}{S1}] ULAM.AI and the UnsolvedMath Contributors, UnsolvedMath, record 30006043 (OWR-14298591-007). Catalogue curation under CC BY 4.0; definitions and English exposition added locally.
\url{https://huggingface.co/datasets/ulamai/UnsolvedMath}.
\item[\hypertarget{source:30006043:2}{S2}] Mathematical Aspects of General Relativity (2024), Oberwolfach Reports, DOI 10.4171/owr/2024/36, pp. 2033--2112. Original report and mathematical context. Incorporated report material remains under the source CC BY-SA 4.0 terms unless otherwise noted.
\url{https://doi.org/10.4171/owr/2024/36}.
\end{itemize}
\label{end:30006043}
\end{document}
